\chapter{Non-Identifiability}
\label{ch:nonidentifiability}

\section{Statement}

The question addressed in this chapter is when it is impossible, not
merely difficult, to recover which of two histories actually occurred
from the evidence either would have left behind. The elementary
direction of this claim is close to definitional; the substantive
direction is showing that the relevant class of observation maps is
genuinely, structurally non-injective, rather than merely
non-injective in unusual or contrived cases.

\begin{definition}[Observation map]
An observation map is a function $\obsmap : \History \to \Observable$
from complete histories to observables, where $\Observable$ is whatever
evidence remains available at a given time to an observer without
access to $\History$ itself.
\end{definition}

\begin{theorem}[Non-Identifiability, elementary direction]
\label{thm:nonident-elementary}
If $H_1 \neq H_2$ and $\obsmap(H_1) = \obsmap(H_2)$, then no function
$f : \Observable \to \History$ satisfies $f(\obsmap(H_1)) = H_1$ and
$f(\obsmap(H_2)) = H_2$ simultaneously.
\end{theorem}

\begin{proof}
Suppose such an $f$ existed. Since $\obsmap(H_1) = \obsmap(H_2)$, write
$o = \obsmap(H_1) = \obsmap(H_2)$. Then $f(o) = f(\obsmap(H_1)) = H_1$
and also $f(o) = f(\obsmap(H_2)) = H_2$, so $H_1 = H_2$, contradicting
the hypothesis $H_1 \neq H_2$.
\end{proof}

This direction is elementary because it is a direct consequence of $f$
being required to be a function, together with the assumed collision
$\obsmap(H_1) = \obsmap(H_2)$; no property of admissibility structure is
needed to prove it. What requires the admissibility apparatus developed
elsewhere in this program is establishing that such collisions are not
pathological exceptions confined to contrived cases, but a structural
feature of the observation maps that actually arise when observation is
constrained to preserve only reachable-set structure.

\begin{theorem}[Non-Identifiability, structural direction]
\label{thm:nonident-structural}
Let $\obsmap$ be any observation map that factors through the
reachable-set structure --- that is, suppose $\obsmap(H) = \psi(\reach{H})$
for some function $\psi$, so that $\obsmap$ records only which
continuations remain admissible from $H$, not $H$ itself. Then whenever
two histories $H_1 \neq H_2$ satisfy $\reach{H_1} = \reach{H_2}$, it
follows that $\obsmap(H_1) = \obsmap(H_2)$, and Theorem
\ref{thm:nonident-elementary} applies. Moreover, this is not a
degenerate case: $\reach{H_1} = \reach{H_2}$ is precisely the condition
under which $H_1$ and $H_2$ lie in the same equivalence class of the
admissibility relation, and any admissibility structure with a
non-trivial quotient --- that is, any admissibility structure in which
more than one history maps to the same reachable set --- guarantees the
existence of such collisions.
\end{theorem}

\begin{proof}
The first claim is immediate: if $\reach{H_1} = \reach{H_2}$ then
$\psi(\reach{H_1}) = \psi(\reach{H_2})$ trivially, since $\psi$ is a
function of the reachable set alone. For the second claim, recall that
$R(\cdot)$ partitions the space of histories into equivalence classes by
definition: two histories lie in the same class exactly when they have
the same reachable set. An admissibility structure has a trivial
quotient only when this partition is the discrete partition, that is,
only when no two distinct histories ever share a reachable set. This is
not the generic case for admissibility structures of the kind developed
elsewhere in this program, where reachable sets are typically determined
by a bounded amount of structure --- the admissible continuations
consistent with what has been committed so far --- rather than by the
full, unbounded historical record. Whenever the reachable set is a
function of bounded information about $H$ while $H$ itself can vary over
an unbounded space, the pigeonhole principle guarantees histories
differing in the unrecorded portion while agreeing in the recorded
portion, hence sharing a reachable set, hence colliding under $\obsmap$.
\end{proof}

\section{Discussion}

Theorem \ref{thm:nonident-structural} is the precise sense in which
Non-Identifiability is not a separate result requiring separate
machinery, but the restorability boundary developed elsewhere in this
program, proven here as a limitative theorem in its own right rather
than argued for informally. The restorability boundary has previously
been discussed in terms of what invariants survive repeated
reconstruction and what information is irrecoverably lost; Theorem
\ref{thm:nonident-structural} identifies exactly the mechanism by which
loss becomes structurally guaranteed rather than merely likely: any
observation map constrained to record reachable-set structure, applied
to a space of histories whose reachable sets are determined by less
information than the histories themselves contain, is non-injective as
a matter of definition, not as a matter of contingent bad luck in any
particular case.

\begin{corollary}
\label{cor:nonident-diagnosis}
Given an observed collision $\obsmap(H_1) = \obsmap(H_2)$ with
$H_1 \neq H_2$, the question of \emph{why} reconstruction fails for this
pair admits a principled decomposition. Either $\reach{H_1} \neq
\reach{H_2}$, in which case the observation map $\obsmap$ has discarded
information that a less lossy observation map could in principle have
retained, and the failure is a failure of the particular observation
procedure; or $\reach{H_1} = \reach{H_2}$, in which case no observation
map factoring through reachable-set structure can distinguish $H_1$ from
$H_2$, and the failure is a failure of recoverability as such, not of
any particular reconstruction procedure.
\end{corollary}

Corollary \ref{cor:nonident-diagnosis} gives the precise form of the
question that motivates this entire research program's treatment of
reconstruction: whether a given case of reconstruction failure reflects
an inference that was performed incorrectly, or evidence that has
structurally ceased to exist. The corollary shows this is not merely a
useful distinction to draw informally but a genuine dichotomy, decidable
in principle by checking whether the colliding histories share a
reachable set.
